\section{Action-aligned acquisition and complete-policy work}\label{sec:action99}
A continuation is acquired to make a decision, not to minimize an approximation statistic in isolation. The retained investment experiment makes the distinction operational: every certificate-oriented fit reduces its validated residual span, whereas aggregate first-proposal rejections increase. We now construct an acquisition target in the units of the admission inequality. The construction uses the same original action set and risk law as the NBO policy account.

\subsection{A computable directional certificate}
Fix a state and date, and suppress their indices. Let $A\subset\mathbb R^k$ be nonempty, compact and convex, let $Q$ be the original action objective, and let $a\in A$ be the actually realized proposal. A verified supporting calculation supplies $L,U\in\mathbb R$, $s\in\mathbb R^k$ and $\mu>0$ such that
\begin{equation}\label{eq:support99}
 Q(a)\le U,\qquad
 Q(b)\ge L+s^{\mathsf T}(b-a)+\frac\mu2\|b-a\|_2^2
 \quad\text{for every }b\in A.
\end{equation}
These inequalities need not identify $s$ as a derivative of $Q$. In particular, an imperfect risk witness or continuation minorant can lower $L$ while retaining a useful supporting slope. Put $w=U-L$ and
\begin{equation}\label{eq:gamma99}
 \Gamma(a)=w-\min_{b\in A}
       \left\{s^{\mathsf T}(b-a)+\frac\mu2\|b-a\|_2^2\right\}.
\end{equation}
A certified lower bound on the displayed minimum gives a valid, possibly larger, $\Gamma$. An unverified optimizer status does not.

\begin{theorem}[Directional certificate for an NBO proposal]\label{thm:action99}
Under \eqref{eq:support99},
\[
 0\le Q(a)-\inf_{b\in A}Q(b)\le\Gamma(a).
\]
Let $\widehat s$ be a supporting slope of the surrogate used to propose $a$, let $N_A(a)=\{n:n^{\mathsf T}(b-a)\le0\ (b\in A)\}$, and suppose
$\|s-\widehat s\|_2\le d$. If a supplied $n\in N_A(a)$ satisfies
$\|\widehat s+n\|_2\le r$, then
\begin{equation}\label{eq:direction99}
 \Gamma(a)\le w+\frac{(d+r)^2}{2\mu}.
\end{equation}
For $A=[0,A_0]$, the exact minimum in \eqref{eq:gamma99} is attained at
$b=\operatorname{proj}_{[0,A_0]}(a-s/\mu)$.
\end{theorem}

The proof, including boundary actions, is in Supplement~\ref{supp:action99}. Equation~\eqref{eq:direction99} separates three trainable or computable quantities: the original-objective enclosure width $w$, the directional mismatch $d$, and the surrogate stationarity residual $r$. The residual span of the continuation is not one of these quantities. A readout can improve that span and worsen the supporting slope at a nearly binding action. Conversely, constant changes in a continuation need not affect any of them after consistent centering.

The supporting calculation is available under the original recursive-risk primitives. Suppose $F(b,\omega)=F(a,\omega)+G(\omega)(b-a)$, the stage cost is $\mu$-strongly convex, and the risk dual has an action-independent feasible witness $q$ on the original innovation space. If
\[
 V(F(b,\omega))\ge \ell(\omega)+v(\omega)^{\mathsf T}G(\omega)(b-a),
\]
then \eqref{eq:support99} holds with
\begin{equation}\label{eq:lift99}
 L=c(a)+\beta\left\{\int\ell\,dq-\psi(q)\right\},\qquad
 s=\nabla c(a)+\beta\int G^{\mathsf T}v\,dq,
\end{equation}
and any valid upper bound $U$ for the original risk objective. Integration and penalty errors must be enclosed in these quantities. The witness must remain feasible for every action being minorized; a witness feasible only at $a$ is insufficient when the risk set itself depends on the action.

In the investment economy, $G=(1/2,1/4)^{\mathsf T}$ at $d=2$, the innovation law and AVaR cap are action independent, and $h''\ge2$. Thus $\mu=2$ gives exactly the all-action quadratic of \eqref{eq:actionlower95}. The added result is the explicit link from an acquired directional support to its admission gap, rather than a new economic model or an assumption that an acquired value support is an exact derivative.

\subsection{Which training loss controls rejection?}
Fix a finite menu $\mathcal H$ of $M$ complete candidate services before validation. Each candidate specifies its proposal, rounding, certificate algorithm and fallback. On a context $x$, its first proposal is accepted when the valid complete-policy certificate $\Gamma_h(x)$ is at most $\tau>0$. For a one-date problem, \eqref{eq:gamma99} is such a certificate. At multiple dates it must additionally include the certified future-policy allowance. Define
\begin{equation}\label{eq:loss99}
 R_h(x)=\boldsymbol 1\{\Gamma_h(x)>\tau\},\qquad
 \ell_h(x)=\min\{1,\Gamma_h(x)/\tau\}.
\end{equation}
Then $R_h\le\ell_h$. Equality at the admission boundary is accepted. A valid fallback completes every rejected context; discarding it would change the service and invalidate its resource comparison.

\begin{theorem}[Prospective rejection and work guarantees]\label{thm:selection99}
Condition on all acquisition information that fixes $\mathcal H$. Let $X_1,\ldots,X_n$ be independent draws from a declared deployment-context law, independent of that information. For $0<\delta<1$, set
\[
 e_n=\sqrt{\frac{\log(2M/\delta)}{2n}}.
\]
An empirical minimizer $\widehat h$ of $R_h$ satisfies, with conditional probability at least $1-\delta$,
\begin{equation}\label{eq:rejection99}
 \mathbb E R_{\widehat h}(X)\le\min_{h\in\mathcal H}\mathbb E R_h(X)+2e_n.
\end{equation}
For an empirical minimizer of $\ell_h$, the same event instead gives
$\mathbb E R_{\widehat h}\le\min_h\mathbb E\ell_h+2e_n$.
If a verified resource bound is $W_h(x)\le C_0+C_1R_h(x)$, with common nonnegative constants, the rejection selector satisfies $\mathbb E W_{\widehat h}\le C_0+C_1(\min_h\mathbb E R_h+2e_n)$, and the surrogate selector has the analogous bound with $\ell_h$. More generally, directly minimizing empirical $W_h/W_{\max}\in[0,1]$ yields
\[
 \mathbb E W_{\widehat h}\le\min_h\mathbb E W_h+2W_{\max}e_n.
\]
All selection and rejected-candidate validation costs are additional one-time acquisition charges.
\end{theorem}

The work statement means $\mathbb E W_{\widehat h}\le C_0+C_1(\min_h\mathbb E R_h+2e_n)$ for the rejection selector, and analogously for $\ell$. A measured maximum is not a verified $W_{\max}$, and a rejection-optimal selector need not minimize runtime when fallback costs differ across contexts. On the finite dyadic domain used by the investment service, a finite fallback can supply a worst-case operation cap; its size, rather than a favorable observed clock, enters this theorem. Conditional concentration allows data-dependent acquisition, but does not allow the menu or certificate rule to be changed after looking at the validation sample.

Theorems~\ref{thm:action99}--\ref{thm:selection99} provide the missing distinction between optimizing a residual certificate and optimizing an action-relevant quantity. They do not retroactively certify the empirical performance of the R97 residual objective, or supply observations for the separately frozen R98 actors. Decision-aware prediction and surrogate calibration are established ideas \citep{ElmachtoubGrigas2022,LiuGrigas2021}. Here the loss is attached to an outward original-law, all-action Bellman certificate and a complete feasible service. The finite-menu concentration argument is classical; the support-to-admission construction specifies what that argument is permitted to select.
